\section{The complete radius and coefficient classification}
\label{sec:phase}

Here and below, an equivalent $a_n\sim C R^{-n}n^{-p}(\log n)^{-q}$
has a strictly positive constant $C$.  This is useful because it proves
that the specified point really is singular, in addition to describing
the coefficients.

\begin{theorem}[Complete feedback phase diagram]\label{thm:phase}
For every $s>0$ and $u>0$, \eqref{eq:model} has a unique holomorphic germ
at zero.  Every coefficient $a_n(s,u)$ is positive.  The following list
gives its exact Taylor radius and coefficient equivalent.  In every
case the positive radius is the unique singularity on the circle of
convergence, and the Taylor series converges absolutely on that circle.

\begin{enumerate}
\item If $0<s\le2$, or if $s>2$ and $u<\Uc(s)$, there is a unique
$t_*\in(0,1)$ such that
\begin{equation}\label{eq:characteristic}
 u^2=2\Li_s(t_*)+\Li_{s-1}(t_*)^2.
\end{equation}
With $w_* =\Li_{s-1}(t_*)$ and $R=t_*\ee^{w_*-u}$,
\begin{equation}\label{eq:interior-law}
 a_n\sim
 \sqrt{\frac{w_*}{2\pi(1+\Li_{s-2}(t_*))}}\,
 R^{-n}n^{-3/2}.
\end{equation}
This is the interior theorem of \cite{RepoFeedback}.

\item Suppose $s>2$ and $u>\Uc(s)$.  Define
\begin{equation}\label{eq:boundary-parameters}
 A=\zeta(s-1),\quad b=\sqrt{u^2-2\zeta(s)},\quad
 \vc=u-b,\quad \Rb=\ee^{-\vc}.
\end{equation}
Then $b>A$, the radius is $\Rb$, and
\begin{equation}\label{eq:supercritical-law}
 \boxed{\displaystyle
 a_n\sim \frac{b^{s-1}}{(b-A)^s}\,\Rb^{-n}n^{-s}.}
\end{equation}
This proves the candidate in the incoming source's first research
question, including every integral $s\ge3$.

\item Suppose $u=\Uc(s)$, $s>2$, so that $b=A$ in
\eqref{eq:boundary-parameters}.  The radius is again $\Rb$.  If
$2<s<3$, put $\alpha=s-1$ and $\rho=1/\alpha$.  Then
\begin{equation}\label{eq:critical-stable}
 a_n\sim
 \frac{(A/\Gamma(-\alpha))^{1/\alpha}}{-\Gamma(-1/\alpha)}
 \,\Rb^{-n}n^{-1-1/\alpha}.
\end{equation}
At $s=3$, with $A=\zeta(2)$,
\begin{equation}\label{eq:critical-log}
 a_n\sim\sqrt{\frac A\pi}\,
 \Rb^{-n}n^{-3/2}(\log n)^{-1/2}.
\end{equation}
If $s>3$, then
\begin{equation}\label{eq:critical-gaussian}
 a_n\sim
 \sqrt{\frac{A}{2\pi(1+\zeta(s-2))}}\,
 \Rb^{-n}n^{-3/2}.
\end{equation}
\end{enumerate}
\end{theorem}

The transition has two independent parameters.  Changing $u$ determines
whether the first obstruction is an interior fold or the endpoint
$t=1$.  Once the endpoint and the fold coincide, changing $s$ determines
the regularity of the endpoint.  In particular, the vanishing limit of
the finite-curvature constant as $s\downarrow3$ does not describe the
$s=3$ law: the logarithm changes the scale.

\subsection{Positivity and the positive real branch}

\begin{lemma}[Positive germ and parametrization]\label{lem:positive}
The unique germ of \eqref{eq:model} has $a_n\ge1/(un^s)>0$.
On the increasing part of the positive real branch it is parametrized by
\begin{equation}\label{eq:parametrization}
 f_u(t)=u-\sqrt{u^2-2\Li_s(t)},\qquad
 v=f_u(t),\qquad r=t\ee^{-f_u(t)}.
\end{equation}
The first stationary point of $r(t)$, when it exists, is exactly
\eqref{eq:characteristic}.
\end{lemma}

\begin{proof}
Write
\[
 \Phi(r,v)=\frac{v^2}{2u}+\frac1u\sum_{n\ge1}\frac{r^n\ee^{nv}}{n^s}.
\]
The implicit function theorem applies to $v=\Phi(r,v)$ at $(0,0)$,
because $\Phi_v(0,0)=0$.  Starting with zero and iterating $\Phi$
formally produces nonnegative coefficients.  If two zero-constant-term
series agree through degree $N-1$, their images agree through degree
$N$, so every coefficient stabilizes after finitely many iterations.
The direct term $r^n/(un^s)$ proves the strict lower bound.

Solving the quadratic at $v=0,t=0$ selects \eqref{eq:parametrization}.
Where its square root is positive,
\begin{equation}\label{eq:radius-derivative}
 \frac{\dd}{\dd t}\log r(t)=\frac1t
 \left(1-\frac{\Li_{s-1}(t)}{\sqrt{u^2-2\Li_s(t)}}\right).
\end{equation}
The function $2\Li_s(t)+\Li_{s-1}(t)^2$ is strictly increasing on
$(0,1)$.  Its limit is infinite for $s\le2$ and equals $\Uc(s)^2$
for $s>2$.  At an interior stationary point the radicand is
$\Li_{s-1}(t)^2>0$, so the stationary point occurs before a possible
zero of the radicand.  These facts prove all assertions.
\end{proof}

\subsection{Global convergence and absence of competing singularities}

We supply the details for $u\ge\Uc(s)$, the previously unproved
range.  The interior proof uses the same argument with $t_*$ in
place of $1$.

For $u\ge\Uc(s)$ the radicand is positive on $[0,1]$, and
\eqref{eq:radius-derivative} is positive for $0<t<1$.
The parametrized solution therefore increases from $(0,0)$ to
$(\Rb,\vc)$.  Fix $0\le r<\Rb$.  The real iterates of $\Phi(r,\cdot)$
starting at zero are bounded above by this solution.  Monotone
convergence, first in the iteration and then in the coefficient
index, shows that the formal series converges at this $r$.
Thus its radius is at least $\Rb$.  Its sum agrees with the
parametrized branch by local uniqueness and continuation.  Taking
$r\uparrow\Rb$ gives
\begin{equation}\label{eq:boundary-sum}
 \sum_{n\ge1}a_n\Rb^n=\vc<\infty.
\end{equation}
Consequently the series converges absolutely and uniformly throughout
the closed disk $|r|\le\Rb$.

Let $|r|=\Rb$, $r\ne\Rb$.  Strict positivity of every $a_n$ gives
$|v(r)|<\vc$: equality in the triangle inequality would force
$r^n/\Rb^n$ to have the same argument for every positive integer $n$.
It follows that
\[
 |r\ee^{v(r)}|<\Rb\ee^{\vc}=1.
\]
The defining equation is analytic at this boundary point, and
\begin{equation}\label{eq:boundary-derivative}
 |\Phi_v(r,v(r))|
 \le\frac{|v(r)|+\Li_{s-1}(|r\ee^{v(r)}|)}u
 <\frac{\vc+A}{u}\le1.
\end{equation}
The implicit function theorem extends the germ across every such
point.  Therefore $\Rb$ is the only possible singularity on the
circle.  The local singular expansions below prove that it is
indeed singular and hence that the radius is exactly $\Rb$.

\subsection{An exact integral normal form}

Define, for $s>2$ and $x$ in a small slit neighborhood of zero,
\begin{equation}\label{eq:DK}
 D_s(x)=\zeta(s)-\Li_s(\ee^{-x}),\qquad
 K_s(x)=A x+\frac{x^2}{2}-D_s(x).
\end{equation}
Since $D_s'(x)=\Li_{s-1}(\ee^{-x})$,
\begin{equation}\label{eq:K-integral}
 \boxed{\displaystyle
 K_s(x)=\int_0^x(x-t)\bigl(1+\Li_{s-2}(\ee^{-t})\bigr)\,\dd t.}
\end{equation}
The integral is along the radial segment on the chosen branch.
It converges at zero for $s>2$, including $2<s\le3$.
Indeed the singular part of its integrand is of order
$t^{s-3}$ or $\log t$, and $s-3>-1$.

Put $x=-\Log t$ and $d=-\Log(r/\Rb)$.  The parametrization gives
the exact equations
\begin{align}
 d&=\Psi_b(x):=x+b-\sqrt{b^2+2D_s(x)},\label{eq:Psi}\\
 v(r)&=\vc+d-x,\label{eq:v-x}\\
 (b+x)d-\frac{d^2}{2}&=(b-A)x+K_s(x).\label{eq:normalform}
\end{align}
The last identity follows by squaring the first one, using the root
that tends to $b$.  Conversely, near zero, that root condition recovers
\eqref{eq:Psi}; no extra local inverse branch is introduced by the
squaring.  At the critical parameter the linear term disappears:
\begin{equation}\label{eq:critical-normalform}
 \boxed{\displaystyle K_s(x)=(A+x)d-\frac{d^2}{2}.}
\end{equation}
This single equation is the starting point of all the convergent
critical charts below.

\subsection{The local singularities}

For nonintegral $s>2$, set $\alpha=s-1$.  Equation
\eqref{eq:jonquiere} gives the convergent identity
\begin{equation}\label{eq:K-noninteger}
 K_s(x)=\Gamma(-\alpha)x^\alpha+x^2 B_s(x),\qquad
 B_s(x)=\frac{1+\zeta(s-2)}2+
 \sum_{k\ge3}\zeta(s-k)\frac{(-1)^k x^{k-2}}{k!}.
\end{equation}
For an integer $m\ge3$ the corresponding resonant formula is
\begin{equation}\label{eq:K-integer}
 K_m(x)=\frac{x^2}{2}
 +\sum_{\substack{k\ge2\\k\ne m-1}}
   \zeta(m-k)\frac{(-x)^k}{k!}
 +\frac{(-x)^{m-1}}{(m-1)!}
   \bigl(H_{m-1}-\Log x\bigr).
\end{equation}
The separate $x^2/2$ remains present when $m=3$.
In that case
\begin{equation}\label{eq:K3}
 K_3(x)=\frac{x^2}{2}\left(\log\frac1x+\frac52\right)
       +\frac{x^3}{12}-\frac{x^4}{288}+O(x^6).
\end{equation}
The absence of the $x^5$ term is $\zeta(-2)=0$.

At $b=A$, \eqref{eq:critical-normalform} now gives
\begin{equation}\label{eq:Psi-critical-leading}
 \Psi_A(x)\sim
 \begin{cases}
  \dfrac{\Gamma(1-s)}A x^{s-1},&2<s<3,\\[4pt]
  \dfrac{x^2}{2A}\log(1/x),&s=3,\\[4pt]
  \dfrac{1+\zeta(s-2)}{2A}x^2,&s>3.
 \end{cases}
\end{equation}
All equivalents hold uniformly on closed subsectors of the slit.
For $b>A$ the leading term is instead
$\Psi_b(x)\sim\kappa x$, where $\kappa=1-A/b>0$.

\begin{lemma}[Sectorial inverse and dented continuation]\label{lem:delta}
Each boundary branch above has a holomorphic inverse $x(d)$ in a
sufficiently small sector $|\arg d|<\pi-\eta$, for any fixed
$0<\eta<\pi$, selected by $x>0$ when $d>0$.  Together with
\eqref{eq:boundary-derivative}, this supplies a dented disk at
$r=\Rb$ of the type required for singularity transfer.
\end{lemma}

\begin{proof}
For $\Psi(x)\sim c x^p$, $c>0$, $p\ge1$, take
$x_0=(d/c)^{1/p}$.  On a small fixed relative circle
$|x/x_0-1|=\epsilon$, the normalized function $\Psi(x)/d-1$
converges uniformly to $(x/x_0)^p-1$, which has one simple zero
inside.  Rouch\'e's theorem gives a unique root near $x_0$.
The differentiated convergent expansions show that the derivative
is nonzero there.  The local implicit inverses glue by uniqueness.
For $\Psi(x)\sim c x^2\log(1/x)$, the same proof uses
$x_0=(2d/(c\log(1/d)))^{1/2}$; the slowly varying logarithm has
ratio tending to one on the relative circle.  These arguments
use the differentiated expansions, rather than differentiating an
arbitrary unspecified $o(1)$.

The variable $d=-\Log(r/\Rb)$ is analytic near $r=\Rb$ and equals
$1-r/\Rb+O((1-r/\Rb)^2)$.  Its slit sector supplies the local
dent.  The rest of the circle has analytic extensions by
\eqref{eq:boundary-derivative}.  A finite cover of its compact
arcs completes a dented disk after its thickness is decreased.
\end{proof}

\subsection{Proof of the coefficient laws}

For nonintegral $s>2$ and $b>A$, regard $x$ and $x^\alpha$ as
independent local coordinates.  In \eqref{eq:Psi}, the coefficient
of its first nonanalytic term is $\Gamma(-\alpha)/b$.
The inverse therefore has first nonanalytic term
\[
 -\frac{\Gamma(-\alpha)}{b\kappa^{\alpha+1}}d^\alpha.
\]
One way to verify the factor is to substitute the analytic inverse
$x_{\rm an}(d)=d/\kappa+O(d^2)$ and divide the perturbation by
$\Psi'_{\rm an}(0)=\kappa$.  Equation \eqref{eq:v-x} changes its
sign.  Thus the singular term of $v$ is
\begin{equation}\label{eq:super-singular}
 \frac{\Gamma(1-s)}{b\kappa^s}d^{s-1}.
\end{equation}
Transfer of $(1-z)^{s-1}$ divides by $\Gamma(1-s)$ and gives
\eqref{eq:supercritical-law}.

The integer case can be checked directly, avoiding any interchange
of a parameter limit and a coefficient asymptotic.  For $s=m\ge3$,
put $j=m-1$.  The singular term of $v$ is
\begin{equation}\label{eq:integer-super-singular}
 \frac{(-1)^{j+1}}{b\,j!\,\kappa^m}d^j\Log d.
\end{equation}
The exact logarithmic transfer equivalent is
\[
 \coef z n (1-z)^j\Log(1-z)
 \sim(-1)^{j+1}j!\,n^{-j-1}.
\]
The signs and factorials cancel, giving the same positive constant
$1/(b\kappa^m)$.

At criticality, inversion of \eqref{eq:Psi-critical-leading} and
\eqref{eq:v-x} yields respectively
\begin{align*}
 v&=\vc-(A/\Gamma(-\alpha))^{1/\alpha}d^{1/\alpha}
       +o(d^{1/\alpha}),&&1<\alpha<2,\\
 v&=\vc-2\sqrt A\,d^{1/2}(\log(1/d))^{-1/2}
       \left(1+O\left(\frac{\log\log(1/d)}{\log(1/d)}\right)\right),
       &&s=3,\\
 v&=\vc-\sqrt{\frac{2A}{1+\zeta(s-2)}}d^{1/2}
       +o(d^{1/2}),&&s>3.
\end{align*}
The term $d$ in \eqref{eq:v-x} is locally analytic and is smaller
than the displayed singular terms.  Apply \cref{lem:delta} and the
power--logarithm transfer theorem \cite{FO}.  Since
$\Gamma(-1/2)=-2\sqrt\pi$, this proves
\eqref{eq:critical-stable}--\eqref{eq:critical-gaussian}.
The local expansions prove genuine singularity and complete the
radius argument.  At an interior point the usual implicit quadratic
expansion has amplitude
$\sqrt{2w_* /(1+\Li_{s-2}(t_*))}$; the same transfer gives
\eqref{eq:interior-law}.  This also reproves the part credited to
the incoming report.

\begin{corollary}[The physical convergence threshold for every order]
\label{cor:physical}
For every $s>0$ there is a unique $t_s\in(0,1)$ satisfying
\begin{equation}\label{eq:physical-t}
 \log t_s+\Li_{s-1}(t_s)=0.
\end{equation}
Put $u_s^2=2\Li_s(t_s)+\Li_{s-1}(t_s)^2$.
The Taylor series of $v$ evaluated at $r=-\ee^{-u}$ converges if and
only if $u\ge u_s$.  At equality it converges absolutely.
\end{corollary}

\begin{proof}
The left side of \eqref{eq:physical-t} is strictly increasing from
$-\infty$ to a positive limit.  In the interior regime,
$\ee^{-u}/R=\ee^{-w_*}/t_*$, so the condition that the evaluation
point lie on or within the convergence circle is precisely
$\log t_*+w_*\ge0$.  Monotonicity of the characteristic equation
proves the stated threshold.  If $s>2$, then $u_s<\Uc(s)$.
For all $u\ge\Uc(s)$ the boundary radius instead satisfies
$\ee^{-u}/\Rb=\ee^{-b}<1$.  Absolute convergence at equality
follows from \cref{thm:phase}.
\end{proof}

The incoming report proves this last statement for $0<s\le2$.
The boundary classification removes that restriction.
